\section{Formal theory}\label{formal-theory}

We prove Theorem~\ref{weighted_ambient_metrics_exist} by iteratively constructing a power series solution to $\widetilde{R}(\widetilde{g}),\widetilde{F_\phi^m}=O(\rho^j)$. This process is only obstructed when $d+m\in 2\mathbb{N}$; in this case the obstruction is at order $O\big(\rho^{\frac{d+m}{2}-1}\big)$.

We first give a necessary condition for a weighted ambient space to be normal.

\begin{Lemma}\label{ambient-metric-base-case}
Let $(\widetilde{\mathcal{G}},\widetilde{g},\widetilde{f},m,\mu)$ be a normal weighted ambient space. Then \begin{equation}\label{standard-normal-form}\widetilde{g}=a\,\mathrm{d}t^2+2b_i\, {\rm d}x^i \,\mathrm{d}t+ 2t\,{\rm d}\rho\, \mathrm{d}t+ t^2 (g_{ij})_\rho,\end{equation} with $a=2\rho+O(\rho^2)$ and $b_i=O(\rho^2)$.
\end{Lemma}

\begin{proof}
Lemma~\ref{straight-metric-lemma} implies that $\widetilde{g}$ has the form of equation~\eqref{standard-normal-form}. Since $\widetilde{g}$ is normal, $a|_{\rho=0}=0$ and $b_i|_{\rho=0}=0$. Denote $\widetilde{R}_{IJ}:=\widetilde{(\ric)}_{IJ}$.
By definition,
\begin{gather*}
\begin{split}
&\widetilde{R}_{IJ}=\frac{1}{2}\widetilde{g}^{KL}\big(\partial^2_{IL}\widetilde{g}_{JK} +\partial^2_{JK}\widetilde{g}_{IL}-\partial^2_{KL}\widetilde{g}_{IJ}-\partial^2_{IJ}\widetilde{g}_{KL}\big)
\\ &\hphantom{\widetilde{R}_{IJ}=}
{}+\widetilde{g}^{KL}\widetilde{g}^{PQ}\big(\widetilde{\Gamma}_{ILP}\widetilde{\Gamma}_{JKQ} -\widetilde{\Gamma}_{IJP}\widetilde{\Gamma}_{KLQ}\big)
-\frac{m}{\widetilde{f}}\big(\partial^2_{IJ}\widetilde{f} -\widetilde{\Gamma}^K_{IJ}\partial_K\widetilde{f}\big).
\end{split}
\end{gather*}
Evaluating at $\rho=0$ yields
\begin{gather*}
\widetilde{R}_{00}= \frac{d+m}{2t^2}(2-\partial_{\rho}a),
\qquad
\widetilde{R}_{0i}=\frac{1}{2t}\partial^2_{i\rho}a -\frac{d+m}{2t}\partial_{\rho}b_i.
\end{gather*}
We conclude that $a=2\rho+O(\rho^2)$ and $b_i=O(\rho^2)$.
\end{proof}

We now take a brief digression to discuss straight weighted ambient spaces. First, notice that a simple choice of certain components of $\widetilde{g}$ makes $\widetilde{R}_{0I}$ vanish.

\begin{Lemma}
\label{straight-zero-theorem}
Let $(\widetilde{\mathcal{G}},\widetilde{g},\widetilde{f},m,\mu)$ be a weighted pre-ambient space. If $\widetilde{g}$ has the form
\begin{equation}
\label{straightmetric}
 \widetilde{g}_{IJ}=\begin{pmatrix} 2\rho&0&t\\0& t^2 g_{\rho}&0\\t&0&0 \end{pmatrix}\!,
\end{equation}
where $g_\rho$ is a one-parameter family of metrics on $M$, then $\widetilde{R}_{0I}=0$.
\end{Lemma}

\begin{proof}This follows by direct computation~(cf.\ \cite[Lemma 3.2]{FeffermanGraham2012}).
\end{proof}

The relevance of Lemma~\ref{straight-zero-theorem} stems from the following equivalent characterizations of straight normal pre-ambient metrics.
\begin{Proposition}
\label{straightproperties}
Let $(\widetilde{\mathcal{G}},\widetilde{g},\widetilde{f},m,\mu)$, $d\geq 3$ and $m<\infty$, be in normal form relative to $(g,f)$. Then the following conditions are equivalent:
\begin{enumerate}[label=$(\roman*)$]\itemsep=0pt
\item $\widetilde{g}_{00}=2\rho$ and $\widetilde{g}_{0i}=0$;

\item for each $p\in \widetilde{\mathcal{G}}$, the dilation orbit $s\mapsto \delta_s p$ is a geodesic for $\widetilde{g}$;

\item $\widetilde{g}(2T,\cdot)={\rm d}(\widetilde{g}(T,T))$;

\item the infinitesimal dilation field $T$ satisfies $\widetilde{\nabla}T=\operatorname{Id}$.
\end{enumerate}
\end{Proposition}

\begin{proof}
The proof is identical to that of \cite[Proposition 3.4]{FeffermanGraham2012}.
\end{proof}

The following result implies that, up to ambient-equivalence, we may restrict our attention to metrics of the form of equation~\eqref{straightmetric}.
\begin{Theorem}
\label{normal-to-straight-and-normal}
A normal weighted ambient space is ambient-equivalent to a straight and normal ambient space.
\end{Theorem}
\begin{proof}
Let $(\widetilde{\mathcal{G}},\widetilde{g},\widetilde{f},m,\mu)$ be a normal weighted ambient space. By Lemma~\ref{ambient-metric-base-case}, we may write
\[
\widetilde{g}=\begin{pmatrix} 2\rho&0&t\\0&t^2 g_\rho& 0\\ t&0&0\end{pmatrix}+\rho^n \begin{pmatrix} a& tb_i& 0\\ tb_j & 0 & 0\\ 0&0&0\end{pmatrix}
\]
for some $n\geq 2$ and some $a=a(x,\rho)$ and $b_i=b_i(x,\rho)$. Direct computation using Lemma~\ref{straight-zero-theorem} yields
\begin{gather*}
t^{2} \widetilde{R}_{00}=n\bigg(n-1-\frac{d+m}{2}\bigg) \rho^{n-1} a+O(\rho^{n}),
\\
t \widetilde{R}_{0 i}=n\bigg(n-1-\frac{d+m}{2}\bigg) \rho^{n-1} b_i+\frac{n}{2} \rho^{n-1} \partial_{i} a+O(\rho^{n}),
\end{gather*}
(cf.\ \cite[equation~(3.11)]{FeffermanGraham2012}).
We conclude that if $d+m\notin 2\mathbb{N}$, then $a,b_i=O(\rho^\infty)$; and if $d+m\in 2\mathbb{N}$, then $n\geq \frac{d+m}{2}$. Therefore $\widetilde{g}$ is ambient-equivalent to a metric of the form of equation~\eqref{straightmetric}. The~conclusion follows from Proposition~\ref{straightproperties}.
\end{proof}

We now iteratively construct a straight and normal weighted ambient space.
Let $n\in\mathbb{N}$ be such that there is a metric
\[
\widetilde{g}_{IJ}^{(n-1)} = \begin{pmatrix}
2\rho & 0 & t \\
0 & t^2 g_{\rho} & 0 \\
t & 0 & 0
\end{pmatrix}
\]
and a function $\widetilde{f}^{(n-1)}=tf_\rho$ such that
\begin{gather*}
\widetilde{R}^{(n-1)}=O\big(\rho^{n-1}\big),
\qquad
\widetilde{F_{\phi}^m}^{(n-1)}=O\big(\rho^{n-1}\big).
\end{gather*}
Note that the existence of $\widetilde{g}_{IJ}^{(1)}$ and $\widetilde{f}^{(1)}$ trivially holds. Let $\cg^{(n)}_{IJ}=\cg^{(n-1)}_{IJ}+\Phi_{IJ}$ and $\widetilde{f}^{(n)}= \widetilde{f}^{(n-1)}+\rho^{n}t\upsilon$, where
\begin{align*}
\Phi_{IJ}&= \rho^n \begin{pmatrix}
0 & 0 & 0 \\
0 & t^2\psi_{ij} & 0 \\
0 & 0 & 0
\end{pmatrix}
\end{align*}
and $\psi_{ij}$, $\upsilon$ depend only on $x$ and $\rho$.
We seek $\psi_{IJ}$ and $\upsilon$ such that
 \begin{gather*}
\widetilde{R}^{(n)}= O(\rho^{n}),
\qquad
\widetilde{F_{\phi}^m}^{(n)} = O(\rho^{n}).
\end{gather*}
Note that the inverse of $\cg^{(n)}$, calculated modulo $O(\rho^n)$, is
\[
\cg^{IJ} = \begin{pmatrix} 0 & 0 & t^{-1}\\ 0 & t^{-2}g^{ij}& 0\\ t^{-1} & 0 & -{2\rho}t^{-2}\end{pmatrix}\!.
\]

The Christoffel symbols for $\widetilde{g}^{(n)}$ modulo $O(\rho^n)$ are (cf.~\cite[equation~(3.16)]{FeffermanGraham2012}).
\begin{gather*}
\widetilde{\Gamma}_{I J }^0=\begin{pmatrix}
0 & 0 & 0 \\
0 & -\frac{1}{2} t\big(\partial_{\rho} g_{i j}+n \rho^{n-1} \psi_{i j}\big) & 0 \\
0 & 0 & 0
\end{pmatrix}\!,
\\
\widetilde{\Gamma}_{I J}^\ell=\begin{pmatrix}
0 & t^{-1} \delta_{j}^{\ell} & 0 \\
t^{-1} \delta_{i}^{\ell} & \Gamma_{i j}^\ell & \frac{1}{2} g^{\ell k} \partial_{\rho} g_{i k}+\frac{n}{2} \rho^{n-1} \psi_{i}^{\ell}\\
0 & \frac{1}{2} g^{\ell k} \partial_{\rho} g_{j k}+\frac{n}{2} \rho^{n-1} \psi_{j}^{\ell} & 0
\end{pmatrix}\!,
\\
\widetilde{\Gamma}_{I J}^\infty=\begin{pmatrix}
0 & 0 & t^{-1} \\
0 & \rho \partial_{\rho} g_{i j}-g_{i j} & 0 \\
t^{-1} & 0 & 0
\end{pmatrix}\!.
\end{gather*}
It follows that~(cf.~\cite[equation~(3.11)]{FeffermanGraham2012})
\begin{subequations}
\begin{gather}
\label{ricij}
\widetilde{R}_{ij}^{(n)}=\widetilde{R}_{ij}^{(n-1)}+n\rho^{n-1} \bigg[\bigg(n-\frac{d+m}{2}\bigg)\psi_{ij}-\frac{1}{2}g^{kl}\psi_{kl}g_{ij}-\frac{m}{f}\upsilon g_{ij}\bigg]+O\big(\rho^{n}\big),
\\
\label{rphi}
\widetilde{F_{\phi}^m}^{(n)}= \widetilde{F_{\phi}^m}^{(n-1)}
+ n\rho^{n-1}f \bigg[\frac{1}{2}fg^{kl}\psi_{kl}+\upsilon (d+2m-2n) \bigg]+O\big(\rho^{n}\big),
\\
\widetilde{R}_{\infty\infty}^{(n)}=\widetilde{R}_{\infty\infty}^{(n-1)}-n(n-1)\rho^{n-2} \bigg[\frac{1}{2}g^{kl}\psi_{kl}+mf^{-1}\upsilon\bigg]+O\big(\rho^{n-1}\big)\label{ricinftyinfty}.
\end{gather}
\end{subequations}

\subsection[Solving widetilde {R}\_\{ij\}, widetilde {F\_\{phi\}\textasciicircum{}m}=O(rho\textasciicircum{}n)]{Solving $\boldsymbol{\widetilde{R}_{ij},\widetilde{F_\phi^m}=O(\rho^n)}$}
\label{section-4.1}

The following theorem specifies the result of recursively determining $(\widetilde{g}^{(n)},\widetilde{f}^{(n)})$ by the requirements $\widetilde{R}_{ij}^{(n)},(\widetilde{F}_\phi^m)^{(n)}=O(\rho^n)$.
\begin{Theorem}
\label{various-possibilities}
Let $(M^d,g,f,m,\mu)$, $d\geq 3$ and $m<\infty$, be a smooth metric measure space. Set $\widetilde{\mathcal{G}}:=\mathbb{R}_+\times M\times (-\epsilon,\epsilon)$.
\begin{enumerate}[label=$(\roman*)$]\itemsep=0pt
\item If $d+m\notin\mathbb{N}$, then there is a straight and normal metric measure structure $(\widetilde{g},\widetilde{f})$ on $\widetilde{\mathcal{G}}$, unique modulo $O(\rho^\infty)$, such that $\widetilde{R}_{ij},\widetilde{F_\phi^m}=O(\rho^\infty)$.

\item If $d+m\in 2\mathbb{N}$, then
\begin{enumerate}\itemsep=0pt
\item[$(a)$] there is a straight and normal metric measure structure $(\widetilde{g},\widetilde{f})$ on $\widetilde{\mathcal{G}}$, unique modulo $O^+\big(\rho^{\frac{d+m}{2}}\big)$, such that $(\widetilde{R}_{ij},\widetilde{F_\phi^m})=O_{ij}^{2,+}\big(\rho^{\frac{d+m}{2}-1}\big)$;
\item[$(b)$] if
\[
\partial_\rho^{(d+m)/2-1}\big|_{\rho=0}\widetilde{R}_{ij}=0,
\]
then there is a straight and normal metric measure structure $(\widetilde{g},\widetilde{f})$ on $\widetilde{\mathcal{G}}$ such that $\widetilde{R}_{ij},\widetilde{F_\phi^m}=O(\rho^{d+m-1})$;

\item[$(c)$] if
\begin{equation}\label{oddconsistency}
\frac{m}{f^2} \widetilde{F_\phi^m}^{(d+m-1)}-g^{ij}\widetilde{R}_{ij}^{(d+m-1)}=O\big(\rho^{d+m}\big),
\end{equation}
then there is a straight and normal metric measure structure $(\widetilde{g},\widetilde{f})$ on $\widetilde{\mathcal{G}}$ such that $\widetilde{R}_{ij},\widetilde{F_\phi^m}=O(\rho^\infty)$.
\end{enumerate}

\item If $d+m\in \mathbb{N}\setminus 2\mathbb{N}$, then
\begin{enumerate}\itemsep=0pt
\item[$(a)$] there is a straight and normal metric measure structure $(\widetilde{g},\widetilde{f})$ on $\widetilde{\mathcal{G}}$, unique modulo $O(\rho^{d+m})$, such that $\widetilde{R}_{ij},\widetilde{F_\phi^m}=O(\rho^{d+m-1})$;
\item[$(b)$] if equation~\eqref{oddconsistency} holds, then there is a straight and normal metric measure structure $(\widetilde{g},\widetilde{f})$ on $\widetilde{\mathcal{G}}$ such that $\widetilde{R}_{ij},\widetilde{F_\phi^m}=O(\rho^\infty)$.
\end{enumerate}
\end{enumerate}
\end{Theorem}

\begin{Remark}
We say that $(\widetilde{g},\widetilde{f})$ is unique modulo $O^+(\rho^k)$ if it is unique modulo $O(\rho^{k})$ and $\frac{1}{2}g^{ij}\widetilde{g}_{ij}+\frac{m}{f}\widetilde{f}$ is unique modulo $O(\rho^{k+1})$.
\end{Remark}

\begin{proof}
On studying equations~\eqref{ricij} and \eqref{rphi}, we observe that if $n\neq \frac{d+m}{2}$, then there is a~unique choice of the trace-free part $\psi_{ij}-\frac{1}{d}g^{kl}\psi_{kl}g_{ij}$ of $\psi_{ij}$ which makes the trace-free part of $\widetilde{R}^{(n)}_{ij}$ vanish modulo $O(\rho^n)$. However, if $n=\frac{d+m}{2}$, then we may not be able to make the trace-free part of $\widetilde{R}_{ij}^{(n)}$ vanish modulo $O(\rho^n)$. Additionally, equations~\eqref{ricij} and \eqref{rphi} imply that
\begin{gather}
g^{ij}\widetilde{R}_{ij}^{(n)}=g^{ij}\widetilde{R}_{ij}^{(n-1)}
+n\rho^{n-1}\frac{(2n-2d-m)}{2}g^{kl}\psi_{kl}-\frac{dmn}{f}\upsilon \rho^{n-1}+O\big(\rho^{n}\big),\nonumber
\\
\widetilde{F_{\phi}^m}^{(n)}= \widetilde{F_{\phi}^m}^{(n-1)}
+ n\rho^{n-1}f \bigg[\frac{1}{2}fg^{kl}\psi_{kl}+\upsilon (d+2m-2n) \bigg]+O\big(\rho^{n}\big).
\label{modified-equations}
\end{gather}
The determinant of the coefficient matrix of $g^{kl}\psi_{kl}$ and $\upsilon$ is
\begin{equation*}(2n-d-m)(n-d-m).\end{equation*}
Therefore if $n\notin\{\frac{d+m}{2},d+m\}$, then there are unique $g^{kl}\psi_{kl}$ and $\upsilon$ such that $g^{ij}\widetilde{R}_{ij}^{(n)},(\widetilde{F_\phi^m})^{(n)}=O(\rho^n)$. However, if $n\in\{\frac{d+m}{2},d+m\}$, then we may not be able to simultaneously solve $g^{ij}\widetilde{R}_{ij}^{(n)}=O(\rho^n)$ and $\widetilde{F_\phi^m}=O(\rho^n)$.

Suppose first that $d+m\notin \mathbb{N}$. Combining the above discussion with Borel's lemma yields a~straight and normal metric measure structure $(\widetilde{g},\widetilde{f})$ such that $\widetilde{R}_{ij},\widetilde{F_\phi^m}=O(\rho^\infty)$.

Suppose next that $d+m\in 2\mathbb{N}$. Set $n=\frac{d+m}{2}$. Equations~\eqref{modified-equations} imply that
\begin{align}
\frac{m}{f^2}\widetilde{F_{\phi}^m}^{(\frac{d+m}{2})}- g^{ij}\widetilde{R}_{ij}^{(\frac{d+m}{2})} ={}&\frac{m}{f^2}\widetilde{F_{\phi}^m}^{(\frac{d+m}{2}-1)}-g^{ij}\widetilde{R}_{ij}^{(\frac{d+m}{2}-1)} \nonumber
\\
&+\frac{(d+m)^2}{2}\rho^{\frac{d+m}{2}-1} \bigg(\frac{1}{2}g^{kl}\psi_{kl}+\frac{m}{f}\upsilon\bigg)+O\big(\rho^{\frac{d+m}{2}}\big).
\label{even-difference}
\end{align}
In particular, there is a unique choice of $\frac{1}{2}g^{kl}\psi_{kl}+\frac{m}{f}\upsilon$ such that the left hand side of equation~\eqref{even-difference} is $O(\rho^{\frac{d+m}{2}})$. Making this choice yields a unique straight and normal metric measure structure $(\widetilde{g},\widetilde{f})$ such that $(\widetilde{R}_{ij},\widetilde{F_\phi^m})=O_{ij}^{2,+}(\rho^{\frac{d+m}{2}-1})$.
Moreover, if the weighted obstruction is zero, then $\widetilde{R}_{ij}^{(n)}=O(\rho^{\frac{d+m}{2}})$. Hence, we may iteratively solve equations~\eqref{ricij} and \eqref{rphi} to obtain a straight and normal metric measure structure $(\widetilde{g},\widetilde{f})$ such that $\widetilde{R}_{ij},\widetilde{F_\phi^m}=O(\rho^{d+m-1})$. The possible obstruction is the same as that discussed in the next case.

Suppose finally that $d+m\in \mathbb{N}\setminus 2\mathbb{N}$. Set $n=d+m$. From the discussion of the first paragraph, we may choose the trace-free part of $\psi_{ij}$ such that $\widetilde{R}_{ij}^{(d+m)}$ is pure trace. Additionally, by equations~\eqref{modified-equations},
\begin{gather*}
%\label{odd-consistency-equations}
g^{ij}\widetilde{R}_{ij}^{(d+m)}=g^{ij}\widetilde{R}_{ij}^{(d+m-1)}+
m(d+m) \rho^{d+m-1}\bigg(\frac{1}{2}g^{kl}\psi_{kl}-\frac{d}{f}\upsilon\bigg)+O\big(\rho^{d+m}\big),
\\
\widetilde{F_{\phi}^m}^{(d+m)}= \widetilde{F_{\phi}^m}^{(d+m-1)}
+f^2(d+m)\rho^{d+m-1} \bigg(\frac{1}{2}g^{kl}\psi_{kl}-\frac{d}{f}\upsilon\bigg)+O\big(\rho^{d+m}\big).
\end{gather*}
Therefore we may choose $g^{ij}\psi_{ij}$ and $\upsilon$ such that $g^{ij}\widetilde{R}_{ij}^{(d+m)},\widetilde{F_{\phi}^m}^{(d+m)}=O(\rho^{d+m})$ if and only if
equation~\eqref{oddconsistency} holds. If equation~\eqref{oddconsistency} holds, then we may continue iteratively improving $(\widetilde{g},\widetilde{f})$ as in the first paragraph. Hence, by Borel's lemma, there is a straight and normal metric measure structure $(\widetilde{g},\widetilde{f})$ such that $\widetilde{R}_{ij},\widetilde{F_\phi^m}=O(\rho^\infty)$.
\end{proof}

\subsection[Solving widetilde {R}\_\{I infty\}=0]{Solving $\boldsymbol{\widetilde{R}_{I\infty}=0}$}
In this subsection, we show that if $(\widetilde{g},\widetilde{f})$ is as in Theorem~\ref{various-possibilities}, then the components $\widetilde{R}_{I\infty}$ can be made to vanish to the appropriate order and that equation~\eqref{oddconsistency} automatically holds. This proves Theorem~\ref{weighted_ambient_metrics_exist}.

\begin{proof}[Proof of Theorem~\ref{weighted_ambient_metrics_exist}]
Let $(\widetilde{g},\widetilde{f})$ be as in Theorem~\ref{various-possibilities}. Equation~\eqref{intro-weighted-bianchi} implies that
\begin{gather}
\widetilde{g}^{JK}\widetilde{\partial}_J\widetilde{R}_{KI} -\widetilde{g}^{JK}\widetilde{\Gamma}_{JK}^Q\widetilde{R}_{QI} -\widetilde{g}^{JK}\widetilde{\Gamma}_{JI}^Q\widetilde{R}_{KQ} -\widetilde{g}^{JK}\widetilde{R}_{IJ}\partial_K \widetilde{\phi}
\!-\frac{1}{2}\partial_I \widetilde{R_{\widetilde{\phi}}^m}-\frac{1}{f^2}\widetilde{F}_\phi^m \partial_I\widetilde{\phi}=0.\!
\label{bianchianalogue}
\end{gather}

Set
\begin{equation*}
n=\begin{cases}
\infty & \text{if}\quad d+m\notin\mathbb{N},\\
\frac{d+m}{2}-1&\text{if}\quad d+m\in 2\mathbb{N},\\
d+m-1& \text{if}\quad d+m\in \mathbb{N}\setminus 2\mathbb{N}.
\end{cases}
\end{equation*}
Taking $I=l$ and $I=\infty$ in equation~\eqref{bianchianalogue} and computing$\;\bmod \;O(\rho^{n})$ yields
\begin{subequations}
\label{bianchieqns}
\begin{gather}
\label{bianchi2}
[d+m-2-2\rho\partial_{\rho}] \widetilde{R}_{\infty l}-\rho g^{ks}g'_{sl}\widetilde{R}_{\infty k}+2\rho \widetilde{R}_{l\infty}\partial_\rho \phi+\rho \partial_l\widetilde{R}_{\infty\infty}=O\big(\rho^{n}\big),
\\
[d+m-2-\rho\partial_{\rho}] \widetilde{R}_{\infty\infty}
-\frac{1}{2}\partial_{\rho}\bigg(g^{ij}\widetilde{R}_{ij}-\frac{m}{f^2}\widetilde{F_\phi^m}\bigg)
+g^{ij}\big(\widetilde{\nabla}_\phi\big)_i\widetilde{R}_{j\infty}\nonumber
\\ \qquad
+\rho\big(2\partial_\rho\phi-g^{ij}g'_{ij}\big)\widetilde{R}_{\infty\infty}
=O(\rho^{n})\label{bianchi3}.
\end{gather}
\end{subequations}

First, suppose that $d+m\notin\mathbb{N}$. We know that $\widetilde{R}_{ij},\widetilde{F_\phi^m}=O(\rho^\infty)$. From equations~\eqref{bianchi2} and~\eqref{bianchi3} we conclude that $\widetilde{R}_{\infty l},\widetilde{R}_{\infty\infty}=O(\rho^\infty)$.

Second, suppose that $d+m\in \mathbb{N}\setminus 2\mathbb{N}$. Equation~\eqref{bianchieqns} implies that $\widetilde{R}_{\infty l},\widetilde{R}_{\infty\infty}=O(\rho^{n})$, and that $g^{ij}\widetilde{R}_{ij}-mf^{-2}\widetilde{F_\phi^m}=O(\rho^{n+1})$. This verifies equation~\eqref{oddconsistency}. Equation~\eqref{ricinftyinfty} gives us the unique value of $\frac{1}{2}g^{kl}\psi_{kl}+\frac{m}{f}\upsilon$ such that $\widetilde{R}_{\infty\infty}=O(\rho^{n})$. Hence, we can now uniquely determine the values of $g^{kl}\psi_{kl}$ and $\upsilon$, and solve $\widetilde{R}_{IJ},\widetilde{F_\phi^m}=O(\rho^\infty)$.

Third, suppose that $d+m\in 2\mathbb{N}$. Equation~\eqref{bianchi2} tells us that $\widetilde{R}_{\infty l}=O(\rho^n)$ and that $\widetilde{R}_{\infty\infty}=O(\rho^{n-1})$. Now recall that $g^{ij}\widetilde{R}_{ij}-mf^{-2}\widetilde{F_\phi^m}=O\big(\rho^{\frac{d+m}{2}}\big)$.
Hence, equation~\eqref{bianchi3} tells us that $\widetilde{R}_{\infty\infty}=O(\rho^n)$.

Finally, since the terms of $(\widetilde{g},\widetilde{f})$ are uniquely determined using the process described above, any two straight and normal weighted ambient structures will be ambient-equivalent to the orders stated in the theorem.
\end{proof}
